\section{Introduction}\label{l:introduction}

An exact linear observation can convey arbitrarily many bits about a
real-valued signal. A finite message computed from several observations
need not convey that much, but its information is not controlled by
the number of observations alone: the prior distribution matters.
If earlier messages have concentrated the signal in a small region,
the next message may exploit that concentration. This paper studies
how to restore a useful geometric spread condition while accounting
for the information revealed in doing so.

The observation model is noiseless Gaussian regression. The unknown
signal is a unit vector \(S\in S^{d-1}\). A row
\(x\sim N(0,I_d)\), independent of the signal and previous rows,
arrives together with its exact label \(\langle x,S\rangle\).
A learner processes these pairs in order. Between observations it
retains one of at most \(2^M\) states; arbitrary computation and fresh
randomness are allowed during a transition. It stops by a deterministic
finite horizon \(T\), and its unit-vector output uses its terminal
state, stopping index, and fresh randomness. The precise measurable
model, including data-independent shared randomness, is given in
Section~\ref{l:model}. Angular accuracy \(\epsilon\) means
\(\arccos\langle S,\widehat S\rangle\le\epsilon\).

Our first information theorem concerns a prior for which the
localization cost can be seen directly. Put \(n=d-1\), take \(Z\) uniformly
in the cube
\[
 [-1/(2\sqrt n),1/(2\sqrt n))^n,
 \qquad S=\phi(Z)=(Z,\sqrt{1-\|Z\|^2}),
\]
and group observations into blocks of \(k=\lfloor n/16\rfloor\)
rows. A message from a block may depend on all its rows and labels
and on preceding messages. If each message has at most
\(\exp(Ad^2)\) possible values, with \(A\) fixed, then the
first \(t\) messages contain at most \(C_A dt\) nats of mutual
information about \(Z\). Theorem~\ref{ld:thm:main} proves a
stronger statement: the same bound holds after each message is
supplemented with a nested dyadic cell containing \(Z\) that restores
the required spread condition. The constant depends on \(A\), and
the theorem gives the exact alphabet condition. The prior is the
image of uniform cube volume, rather than uniform area on the
spherical patch. This theorem serves as a complete first example; each later
route establishes its own information estimate from its stated
assumptions.

\subsection{Why localization can pay for itself}

Suppose the current posterior is supported in a dyadic parent cell.
It is sufficiently spread for our projection estimate when each
descendant at relative depth \(j\) has probability at most
\(B2^{-nj/2}\). A fresh block then carries at most
\(Cd+\log B+(2/m+e^{-ck})\log N\) information in a message
with \(N\) possible values, where \(m=\lfloor n/8\rfloor\).
The proof first inserts arbitrary \(L^2\) weights into the
posterior, bounds the exact projected density in \(L^m\), and
then uses duality to control the message likelihoods.

After the message, we reveal the deepest cell on the sampled point's
dyadic chain whose posterior mass is at least \(2^{-nj/2}\).
At depth \(j\) there are at most \(2^{nj/2}\) such cells.
Naming the selected cell therefore costs about \(nj/2\) bits,
plus the entropy of its depth. In contrast, confinement to any
depth-\(j\) cell certifies \(nj\) bits relative to the initial
uniform cube. This difference pays for the extra revelations.
The normalization is by the probability of selecting the cell,
not merely the probability of lying in it. Controlling that
selection probability is what also pays for the next value of
\(\log B\).

Section~\ref{l:dyadic} completes this argument and its finite-state
application. The next eight arguments change the localization or
geometric input, with the distinct estimates in Table~\ref{l:routes}.
The final two routes instead correct the information itself. Each
route has its own hypotheses and parameters.

\newpage
\begin{table}[!ht]
\centering
\small
\begin{tabular}{@{}p{0.28\linewidth}p{0.66\linewidth}@{}}
\toprule
Construction & Estimate or information retained \\
\midrule
\multicolumn{2}{@{}l}{\emph{The cube-chart prior}} \\
\raggedright Maximum score\newline
(Section~\ref{l:maxscore:section}) &
A finite maximizing-cell selection, with a one-block information
bound charging half the concentration slack. \\
\addlinespace
\raggedright Finite entropy\newline
(Section~\ref{l:finite-entropy:section}) &
Regularization of a possibly unbounded density of finite relative
entropy, with finite expected termination and the complete stopping
transcript retained. \\
\addlinespace
\raggedright Bounded-density split\newline
(Section~\ref{l:bounded-split:section}) &
A finite partition and a depth--entropy balance; two independent
block copies at the same signal control a squared state likelihood. \\
\addlinespace
\raggedright Likelihood truncation\newline
(Section~\ref{sec:l-likelihood-truncation}) &
A finite-horizon success bound from explicitly bounded exceptional
mass and a pointwise density bound on the surviving histories. \\
\midrule
\multicolumn{2}{@{}l}{\emph{The uniform spherical prior}} \\
\raggedright Grid refinement\newline
(Section~\ref{l:grid:section}) &
Unconditional information in the full history, including the entire
selected grid-cell label after each block. \\
\addlinespace
\raggedright Side projection\newline
(Section~\ref{l:side:section}) &
Information in the full state history conditional on one fixed
Gaussian side observation; predicting adaptive cells bounds the
conditional entropy of their labels. \\
\addlinespace
\raggedright Separated tuples\newline
(Section~\ref{l:sep:section}) &
A terminating radius refinement and an equal-label tuple of controlled
affine volume; a sliced-entropy recurrence for at most \(d\) blocks. \\
\addlinespace
\raggedright Predecessor spans\newline
(Section~\ref{l:span:section}) &
A selected posterior draw and a discrete smaller-ball code, paid by
radius reduction, for possibly unbounded finite-entropy posteriors. \\
\bottomrule
\end{tabular}
\caption{Eight localization estimates following the dyadic proof.
The group labels indicate the priors in the learner applications;
the local estimates apply to the conditional laws specified in each
section. The side-observation matrix is drawn independently of the
learner's observations and the signal. All auxiliary cells,
projections, and posterior draws belong to the analysis;
the learner retains only its prescribed finite state.}
\label{l:routes}
\end{table}

For the first information correction, let \(\nu\) be a bounded-density
probability on the sphere and put \(p=2\lfloor d/16\rfloor\). Define
\[
 J_\nu(s)=\int\|u-s\|^{-p}\,d\nu(u),\qquad
 \Phi(\nu)=D(\nu\|\sigma)-\int\log J_\nu(s)\,d\nu(s),
\]
where \(\sigma\) is uniform spherical probability and \(D\) denotes
relative entropy. Section~\ref{lk:section} proves that \(\Phi\) is
convex, controls a fixed fraction of relative entropy up to \(O(d)\),
and grows by at most \(O(d)\) per block in its stated width range.
Convexity permits averaging away the entering state, so the iteration
follows the current padded state, including its terminal label and
stopping index. Section~\ref{lc:cell} instead keeps the complete
boundary-state history and subtracts a correction built from the
largest cell mass weighted by scale. Its geometric input bounds an
inverse simplex volume under local ball bounds with no support
restriction. The two corrections use different conditioning variables;
neither is identified with the entropy cost of the dyadic revelations.

For \(M(d)=o(d^2)\), the learner applications give
\[
 T(d)\ge c d\log(1/\epsilon(d)),\qquad
 0<\epsilon(d)\le1/10,
\]
with an absolute \(c>0\) and a dimension threshold that may depend
on the memory sequence. The cube arguments assume success at least
\(2/3\) for every signal, which supplies both their cube prior and
the full-sphere prior used at constant accuracy. The spherical
arguments assume average success at least \(2/3\) for a uniform
signal. Each application accounts for the stopping index and the
constant-accuracy endpoint. For the controlled-horizon estimate,
output counting first puts any putative shorter run within the
required range.

\subsection{Context and companion inputs}

Memory constraints distinguish what a learner can extract from a sample
from what it can retain for later samples. Steinhardt and Duchi
\cite{SteinhardtDuchi2015} established memory--sample tradeoffs for noisy
sparse regression. Steinhardt, Valiant, and Wager
\cite{SteinhardtValiantWager2016} formulated a parity-learning conjecture;
Raz \cite{Raz2016} proved a quadratic-memory versus exponential-sample
separation and subsequently extended the method to broader finite
learning problems \cite{Raz2017}. For real-valued linear prediction,
Dagan, Kur, and Shamir \cite[Section~3.3, Theorem~10]{DaganKurShamir2019}
proved a quadratic-memory lower bound for approximately solving a
consistent linear system whose rows arrive once in random order.
Their guarantee concerns residual error on that fixed system, whereas
here the signal is fixed before independent Gaussian samples are drawn.

Sharan, Sidford, and Valiant \cite{SSV2019} studied the statistical
Gaussian-regression problem with small positive observation noise.
Their Theorem~1 gives a sample lower bound of \(\Omega(d\log r)\)
for Euclidean error \(\epsilon=d^{-r}\), memory at most \(d^2/4\)
bits, and \(r\le O(d/\log d)\). Their Section~7, Lemma~12 expands
projection moments into independent signal copies and controls them
by successive orthogonalization. This is a close methodological
predecessor of the projection arguments below. We prove the exact
multirow density estimates, their limiting passages, and the costs
of restoring geometric spread locally. The separate condition-number
conjecture in their Section~1.1 concerns exact labels and
ill-conditioned Gaussian covariance at constant error; the covariance
here is the identity and the lower bound tracks the required accuracy.

The information layer uses relative entropy on standard Borel spaces.
Austin \cite{Austin2020} gives the conditional chain rules and
product decompositions that make these identities meaningful without
a finite signal alphabet. The variational inequality used to charge
selected laws is recorded in \cite{DupuisMao2022}; at each use we
prove the necessary integrability or truncation step for the
unbounded logarithm. The tuple arguments use total correlation in
Watanabe's terminology \cite{Watanabe1960}. Spherical slicing is
related to Federer's coarea formula \cite{Federer1959}, and
Section~\ref{l:sep:section} gives the specific surface-area
calculation. Its subset-entropy comparison credits the entropy
covering argument of Chung, Graham, Frankl, and Shearer
\cite{ChungGrahamFranklShearer1986}; the conditional
differential-entropy argument is proved here. The predecessor-span
code uses Kraft's inequality \cite{Kraft1949}, with the countable
form in \cite[Theorem~2.4.2]{Duchi2023}, followed by a local
conditional entropy calculation.

Only two long analytic inputs are taken from companion articles.
The finite-measure identity of \cite[Theorem~3.3]{OpenAIReplica2026}
supplies the exact equal-label reference in
Section~\ref{l:span:section}. We apply its uniform spherical case
at two separate row counts and prove the finite-entropy extension
locally. The mixed projection moment of
\cite[Theorem~3.1]{OpenAIRow2026} supplies the positive-scale
estimate in Section~\ref{lk:section}. We verify its two local-mass
bounds after reweighting, change the actual-row law while the scale
is positive, and then take the exact-label limit under the actual
law. The remaining localization arguments are proved in this paper.

\subsection{Reading the arguments}

Sections~\ref{l:model} and~\ref{l:dyadic} give a complete first proof:
the model, measurable-rule reduction, information conventions, and
endpoint, followed by dyadic localization and its learner application.
The four other cube routes require this common model and the cube
geometry of Lemma~\ref{ld:cube-geometry}, but not the preceding cube
proofs. The spherical routes can be read after the common model.
In particular, the side-projection argument uses the statement of
Corollary~\ref{l:grid:precision}, not the grid proof, and the two tuple
arguments are independent of each other.

Only the predecessor-span and kernel routes use the companion inputs
described above, stated locally in
Proposition~\ref{l:span:equal-label-assumption} and
Theorem~\ref{lk:mixed-input}. The other geometric and localization
estimates are proved here. Each route fixes its parameters at its
outset, and its final application states the prior and stopping
conventions used for the sample bound.
